\documentclass[11pt]{article}

\usepackage[margin=1.05in]{geometry}
\usepackage{amsmath,amssymb,amsthm,mathtools}
\usepackage{microtype}
\usepackage[hidelinks]{hyperref}
\usepackage{aliascnt}
\usepackage[nameinlink,noabbrev]{cleveref}
\usepackage{enumitem}

\newtheorem{theorem}{Theorem}[section]

\newaliascnt{proposition}{theorem}
\newtheorem{proposition}[proposition]{Proposition}
\aliascntresetthe{proposition}

\newaliascnt{lemma}{theorem}
\newtheorem{lemma}[lemma]{Lemma}
\aliascntresetthe{lemma}

\newaliascnt{corollary}{theorem}
\newtheorem{corollary}[corollary]{Corollary}
\aliascntresetthe{corollary}

\newaliascnt{remark}{theorem}
\newtheorem{remark}[remark]{Remark}
\aliascntresetthe{remark}

\newaliascnt{definition}{theorem}
\newtheorem{definition}[definition]{Definition}
\aliascntresetthe{definition}

\newaliascnt{example}{theorem}
\newtheorem{example}[example]{Example}
\aliascntresetthe{example}

\numberwithin{equation}{section}

\DeclareMathOperator{\Hom}{Hom}
\DeclareMathOperator{\homnum}{hom}
\newcommand{\bbE}{\mathbb E}
\newcommand{\cH}{\mathcal H}
\newcommand{\cF}{\mathcal F}
\newcommand{\cP}{\mathcal P}
\newcommand{\dd}{\,\mathrm d}
\newcommand{\K}{K}
\newcommand{\PhiH}{\Phi_H}
\newcommand{\PsiH}{\Psi_H}
\newcommand{\N}{\mathbb N}

\title{An Entropy Extension of the Chromatic-Polynomial\\
Lower Bound for Chordal Graphs}
\author{}
\date{}

\begin{document}
\maketitle

\begin{abstract}
Let $H$ be a connected chordal graph of clique number $r\geq 3$, and let
$W$ be a graphon of edge density $p$.  The known entropy proof of
\[
 t(H,W)\geq (1-p)^{v(H)}\chi_H\!\left(\frac{1}{1-p}\right)
\]
assumes that all maximal cliques of $H$ have size $r$.  We give a strictly
more general sufficient condition.  If $C_1,\dots,C_m$ are the maximal
cliques of $H$, $T$ is a clique tree, $r_i=|C_i|$, and
$s_e=|C_i\cap C_j|$ for $e=ij\in E(T)$, define
\[
 b_j:=\#\{i:r_i\geq j\}-\#\{e:s_e\geq j\},
 \qquad 1\leq j\leq r.
\]
We prove that the chromatic-polynomial lower bound holds throughout the
usual Tur\'an range $p\geq 1-1/(r-1)$ whenever
$b_1\leq b_2\leq\cdots\leq b_r$.  The proof replaces the separate uniform
clique laws by marginals of one uniform $K_r$-homomorphism law, and then
uses a suffix-majorization consequence of the Moon--Moser clique-ratio
inequalities.  We also give an explicit, generally weaker lower bound for
every connected chordal graph by pooling the decreasing blocks of the
sequence $(b_j)$.
\end{abstract}

\section{Definitions and statements}

A \emph{graphon} is a symmetric measurable function
$W:[0,1]^2\to[0,1]$.  For a finite graph $F$, its homomorphism density in
$W$ is
\[
 t(F,W)
 :=
 \int_{[0,1]^{V(F)}}
 \prod_{uv\in E(F)}W(x_u,x_v)
 \prod_{u\in V(F)}\dd x_u.
\]
We write
\[
 p:=t(\K_2,W)
\]
for the edge density of $W$.

For $0\leq p<1$, define
\[
 \PhiH(p)
 :=
 (1-p)^{v(H)}\chi_H\!\left(\frac{1}{1-p}\right),\label{eq:def-Phi}
\]
where $\chi_H$ denotes the chromatic polynomial of $H$.  We set
$\PhiH(1):=1$, which is the continuous extension of
\eqref{eq:def-Phi} to $p=1$.

A graph $H$ is \emph{chordal} if every cycle of length at least four has a
chord.  Equivalently, its maximal cliques admit a \emph{clique tree}: a
tree $T$ whose vertices are the maximal cliques $C_1,\dots,C_m$ of $H$,
such that, for every $v\in V(H)$, the set
\[
 \{i:v\in C_i\}
\]
induces a connected subtree of $T$.  Throughout the note, $H$ is connected
and chordal.  We write
\[
 r_i:=|C_i|,
 \qquad
 r:=\omega(H)=\max_i r_i,
\]
and, for an edge $e=ij\in E(T)$, we write
\[
 S_e:=C_i\cap C_j,
 \qquad
 s_e:=|S_e|.
\]
For $1\leq j\leq r$, define
\begin{equation}
 b_j
 :=
 \#\{i:r_i\geq j\}
 -
 \#\{e\in E(T):s_e\geq j\}.\label{eq:def-b}
\end{equation}
We shall prove in \cref{sec:chordal} that the integers $b_j$ are positive
and independent of the chosen clique tree.

For $1\leq j\leq r$ and $p\in[0,1]$, set
\begin{equation}
 c_j(p):=1-(j-1)(1-p).\label{eq:def-c}
\end{equation}
Thus $c_1(p)=1$, $c_2(p)=p$, and
\[
 A_s(p):=\prod_{j=1}^s c_j(p)
 =\prod_{a=0}^{s-1}\bigl(1-a(1-p)\bigr).
\]

Our main result is the following.

\begin{theorem}[Monotone exponent condition]\label{thm:main}
Let $H$ be a connected chordal graph with clique number $r\geq3$, and let
$b_1,\dots,b_r$ be defined by \eqref{eq:def-b}.  Suppose that
\begin{equation}
 b_1\leq b_2\leq\cdots\leq b_r.\label{eq:b-monotone}
\end{equation}
Then every graphon $W$ with edge density
\[
 p=t(\K_2,W)\geq 1-\frac{1}{r-1}
\]
satisfies
\begin{equation}
 t(H,W)
 \geq
 \PhiH(p)
 =
 \prod_{j=1}^r c_j(p)^{b_j}.\label{eq:main-bound}
\end{equation}
For $p<1$, the product on the right is equal to
$(1-p)^{v(H)}\chi_H(1/(1-p))$.
\end{theorem}

The pure chordal theorem is an immediate special case.

\begin{corollary}[Pure chordal graphs]\label{cor:pure}
Suppose every maximal clique of $H$ has the same size $r\geq3$.  Then
\eqref{eq:main-bound} holds for every graphon of edge density
$p\geq1-1/(r-1)$.
\end{corollary}

The condition in \cref{thm:main} is strictly weaker than purity; see
\cref{sec:structure}.  In \cref{sec:all-chordal}, we also derive an
explicit lower bound for every connected chordal graph, without assuming
\eqref{eq:b-monotone}.

\begin{remark}[Logical dependence and status]\label{rem:status}
All chordal and entropy arguments are proved below.  The only substantive
external inequality is the Moon--Moser clique-ratio inequality, stated in
\cref{thm:moon-moser}.  No claim concerning priority in the literature is
made for \cref{thm:main} or \cref{thm:all-chordal}.
\end{remark}

\section{Clique trees and the chromatic polynomial}\label{sec:chordal}

We first record the clique-tree identities needed later.

\begin{lemma}[Clique-tree factorization]\label{lem:factorization}
Let $H$ be a connected chordal graph with maximal cliques
$C_1,\dots,C_m$ and clique tree $T$.  Then
\begin{equation}
 v(H)
 =
 \sum_{i=1}^m r_i
 -
 \sum_{e\in E(T)}s_e,\label{eq:vertex-inclusion-exclusion}
\end{equation}
and
\begin{equation}
 \chi_H(x)
 =
 \frac{\prod_{i=1}^m (x)_{r_i}}
      {\prod_{e\in E(T)}(x)_{s_e}},\label{eq:chordal-chi-factorization}
\end{equation}
where
\[
 (x)_k:=x(x-1)\cdots(x-k+1),
 \qquad
 (x)_0:=1.
\]
\end{lemma}

\begin{proof}
For \eqref{eq:vertex-inclusion-exclusion}, fix $v\in V(H)$ and let
\[
 I_v:=\{i:v\in C_i\}.
\]
By the running-intersection property, $T[I_v]$ is a nonempty tree.  Hence
it has $|I_v|-1$ edges.  The contribution of $v$ to the right-hand side
of \eqref{eq:vertex-inclusion-exclusion} is therefore
\[
 |I_v|-\bigl(|I_v|-1\bigr)=1.
\]
Summing over all $v\in V(H)$ proves
\eqref{eq:vertex-inclusion-exclusion}.

We prove \eqref{eq:chordal-chi-factorization} by induction on the number
$m$ of maximal cliques.  If $m=1$, then $H=K_{r_1}$ and
$\chi_H(x)=(x)_{r_1}$.

Suppose $m\geq2$.  Choose a leaf clique $C_m$ of $T$, let $C_k$ be its
unique neighbour in $T$, and put
\[
 S:=C_m\cap C_k,
 \qquad
 s:=|S|,
 \qquad
 R:=C_m\setminus S.
\]
Every vertex of $R$ belongs to no maximal clique other than $C_m$.
Indeed, if $v\in R$ also belonged to a maximal clique $C_i\neq C_m$,
then every clique on the path from $C_m$ to $C_i$ would contain $v$.
In particular, $C_k$ would contain $v$, contradicting $v\notin S$.
Since every edge of $H$ is contained in a maximal clique, it follows that
every neighbour of a vertex in $R$ lies in $C_m$.

Consequently, after a proper coloring of $H-R$ has been fixed, the
vertices of $S$ have $s$ pairwise distinct colors, and the vertices of
$R$ form a clique of size $r_m-s$ adjacent to all of $S$ and to no
additional vertex of $H-R$.  The number of ways to extend the coloring
to $R$ is therefore
\[
 (x-s)_{r_m-s}
 =
 \frac{(x)_{r_m}}{(x)_s}.
\]
Deleting the leaf $C_m$ from $T$ gives a clique tree for $H-R$.  Hence,
for every sufficiently large positive integer $x$, induction gives
\[
 \chi_H(x)
 =
 \frac{(x)_{r_m}}{(x)_s}
 \frac{\prod_{i=1}^{m-1}(x)_{r_i}}
      {\prod_{e\in E(T-C_m)}(x)_{s_e}}
 =
 \frac{\prod_{i=1}^m(x)_{r_i}}
      {\prod_{e\in E(T)}(x)_{s_e}}.
\]
Both sides are polynomials in $x$, so equality for all sufficiently large
positive integers implies the polynomial identity
\eqref{eq:chordal-chi-factorization}.
\end{proof}

\begin{lemma}[The integers $b_j$]\label{lem:b-properties}
For $1\leq j\leq r$, let $b_j$ be given by \eqref{eq:def-b}.  Then:
\begin{enumerate}[label=\textup{(\roman*)}]
\item
\begin{equation}
 \chi_H(x)
 =
 \prod_{j=1}^r (x-j+1)^{b_j};\label{eq:chi-b-factorization}
\end{equation}
\item
\begin{equation}
 v(H)=\sum_{j=1}^r b_j;\label{eq:v-sum-b}
\end{equation}
\item each $b_j$ is a positive integer, $b_1=1$, and the sequence
$(b_1,\dots,b_r)$ is independent of the chosen clique tree;
\item for every $p\in[0,1]$,
\begin{equation}
 \PhiH(p)=\prod_{j=1}^r c_j(p)^{b_j},\label{eq:Phi-product}
\end{equation}
where the right-hand side is understood as the continuous definition at
$p=1$.
\end{enumerate}
\end{lemma}

\begin{proof}
Since
\[
 (x)_s=\prod_{j=1}^s(x-j+1),
\]
substitution into \eqref{eq:chordal-chi-factorization} shows that the
exponent of $x-j+1$ is exactly
\[
 \#\{i:r_i\geq j\}-\#\{e:s_e\geq j\}=b_j,
\]
which proves \eqref{eq:chi-b-factorization}.  Similarly,
\[
 \sum_{j=1}^r b_j
 =
 \sum_i\sum_{j=1}^{r_i}1
 -
 \sum_e\sum_{j=1}^{s_e}1
 =
 \sum_i r_i-\sum_e s_e
 =v(H),
\]
proving \eqref{eq:v-sum-b}.

For positivity, let $T_j$ be the forest with
\[
 V(T_j):=\{i:r_i\geq j\},
 \qquad
 E(T_j):=\{e\in E(T):s_e\geq j\}.
\]
If $e=ik$ satisfies $s_e\geq j$, then both $r_i$ and $r_k$ are at least
$j$, so $T_j$ is indeed a subgraph of the tree $T$.  Since $j\leq r$,
its vertex set is nonempty.  Therefore
\[
 b_j=|V(T_j)|-|E(T_j)|
\]
is the number of connected components of the nonempty forest $T_j$.
In particular, $b_j\geq1$.  We also have $b_1=1$.  Indeed, every
separator of a clique tree of a connected graph is nonempty: if an edge
of $T$ had empty separator, deleting that tree edge would partition the
maximal cliques into two families with disjoint vertex unions, and no edge
of $H$ could join the two unions because every graph edge lies in a
maximal clique.  This would contradict connectedness of $H$.  Hence
$T_1=T$, and therefore $b_1=|V(T)|-|E(T)|=1$.

Independence of the clique tree follows also from
\eqref{eq:chi-b-factorization}, because the multiplicity of each linear
factor $x-j+1$ in $\chi_H(x)$ is intrinsic to $H$.

Finally, let $q:=1-p$.  For $p<1$, equations
\eqref{eq:def-Phi}, \eqref{eq:chi-b-factorization}, and
\eqref{eq:v-sum-b} give
\[
 \begin{aligned}
 \PhiH(p)
 &=q^{v(H)}
   \prod_{j=1}^r\left(\frac{1}{q}-j+1\right)^{b_j}\\
 &=q^{v(H)-\sum_jb_j}
   \prod_{j=1}^r\bigl(1-(j-1)q\bigr)^{b_j}\\
 &=\prod_{j=1}^r c_j(p)^{b_j}.
 \end{aligned}
\]
At $p=1$, every factor on the right equals one, which is the continuous
extension of the same expression.
\end{proof}

\section{Finite graphs, entropy, and the Moon--Moser input}

For a finite graph $G$ on $n$ vertices, write
\[
 \homnum(F,G):=|\Hom(F,G)|,
 \qquad
 t(F,G):=\frac{\homnum(F,G)}{n^{v(F)}}.
\]
All graphs are finite, simple, and loopless.  We use natural logarithms
throughout.

We repeatedly use the following elementary entropy facts.  If a finite
random variable $Z$ is supported on a set of cardinality at most $M$,
then
\begin{equation}
 H(Z)\leq\log M.\label{eq:support-entropy}
\end{equation}
For jointly distributed finite random variables,
\begin{equation}
 H(Z_1,\dots,Z_k)
 =
 \sum_{j=1}^k H(Z_j\mid Z_1,\dots,Z_{j-1}),\label{eq:entropy-chain}
\end{equation}
and conditioning cannot increase entropy.

\subsection{The junction-tree entropy lemma}

\begin{lemma}[Junction-tree entropy gluing]\label{lem:junction-entropy}
Let $H$ be a connected chordal graph with maximal cliques
$C_1,\dots,C_m$ and clique tree $T$, and let $G$ be a finite graph.  For
each $i$, let $\mu_i$ be a probability measure on
$\Hom(H[C_i],G)$.  Suppose that, for every edge $e=ij\in E(T)$, the
marginals of $\mu_i$ and $\mu_j$ on
$S_e=C_i\cap C_j$ agree; denote their common marginal by $\nu_e$.
Then
\begin{equation}
 \log\homnum(H,G)
 \geq
 \sum_{i=1}^m H(\mu_i)
 -
 \sum_{e\in E(T)}H(\nu_e).\label{eq:junction-entropy}
\end{equation}
\end{lemma}

\begin{proof}
Root $T$ at a clique $C_\rho$.  We construct a random assignment
$X:V(H)\to V(G)$ recursively.  First sample $X|_{C_\rho}$ according to
$\mu_\rho$.  Suppose that $C_i$ is a non-root clique, its parent is
$C_{\pi(i)}$, and
\[
 S_i:=C_i\cap C_{\pi(i)}.
\]
When the values of $X$ on all earlier cliques have been sampled, the
values on $S_i$ are already determined.  Conditional on these values,
sample the coordinates in $C_i\setminus S_i$ using the conditional law
of $\mu_i$ given its restriction to $S_i$.  Conditional laws on
separator states of probability zero may be defined arbitrarily; the
construction never reaches such states because the separator marginal
of the parent law equals that of $\mu_i$.

The running-intersection property guarantees that whenever a vertex has
already been assigned, every later occurrence of that vertex lies in the
relevant separator.  Thus the construction is consistent.  Moreover,
each edge of $H$ lies in a maximal clique, and the assignment on that
clique is almost surely a homomorphism into $G$.  Hence the resulting
assignment $X$ lies in $\Hom(H,G)$ almost surely.

By induction away from the root, the marginal law of $X|_{C_i}$ is
$\mu_i$ for every $i$.  Applying the entropy chain rule in the rooted
clique order gives
\[
 \begin{aligned}
 H(X)
 &=H(X|_{C_\rho})
   +\sum_{i\neq\rho}
     H\bigl(X|_{C_i\setminus S_i}\mid X|_{S_i}\bigr)\\
 &=H(\mu_\rho)
   +\sum_{i\neq\rho}\bigl(H(\mu_i)-H(\nu_{i\pi(i)})\bigr)\\
 &=\sum_{i=1}^mH(\mu_i)-\sum_{e\in E(T)}H(\nu_e).
 \end{aligned}
\]
Since $X$ is supported on the set $\Hom(H,G)$,
\eqref{eq:support-entropy} yields
$H(X)\leq\log\homnum(H,G)$.  This proves
\eqref{eq:junction-entropy}.
\end{proof}

\subsection{The clique-ratio inequality}

The only external extremal inequality used in the proof is the following
form of the Moon--Moser theorem.

\begin{theorem}[Moon--Moser clique ratios]\label{thm:moon-moser}
Let $G$ be a finite graph, write
\[
 p:=t(\K_2,G),
 \qquad
 t_j:=t(\K_j,G),
 \qquad
 t_0=t_1:=1.
\]
If $j\geq2$ and $c_j(p)=1-(j-1)(1-p)\geq0$, then
\begin{equation}
 t_j\geq c_j(p)t_{j-1}.\label{eq:moon-moser}
\end{equation}
Consequently, whenever $0\leq k<r$ and all the factors
$c_{k+1}(p),\dots,c_r(p)$ are nonnegative,
\begin{equation}
 \frac{t_r}{t_k}
 \geq
 \prod_{j=k+1}^r c_j(p)
 =
 \frac{A_r(p)}{A_k(p)},\label{eq:moon-moser-iterated}
\end{equation}
provided the quotient is interpreted only in the case $t_k>0$.
\end{theorem}

In the strict range
\[
 p>1-\frac{1}{r-1},\label{eq:strict-range}
\]
all $c_j(p)$, $1\leq j\leq r$, are positive, and
\eqref{eq:moon-moser} implies $t_r\geq A_r(p)>0$.  In particular,
$t_k>0$ for every $k\leq r$, and all quotients in
\eqref{eq:moon-moser-iterated} are well defined.

\subsection{Passage from finite graphs to graphons}

For completeness, we record the approximation statement used at the end
of the proof.

\begin{lemma}[Finite graph approximation]\label{lem:finite-approximation}
Let $W$ be a graphon and let $F_1,\dots,F_q$ be finitely many finite
graphs.  There exists a sequence of finite simple graphs $(G_n)$ such
that
\[
 t(F_a,G_n)\longrightarrow t(F_a,W)
 \qquad
 (1\leq a\leq q).
\]
\end{lemma}

\begin{proof}
For each $n$, sample independent uniform points
$X_1,\dots,X_n\in[0,1]$.  Conditional on these points, join each pair
$1\leq i<j\leq n$ independently with probability $W(X_i,X_j)$.  Let
$\mathbb G(n,W)$ denote the resulting random simple graph.

Fix a finite graph $F$ with $v=v(F)$.  Expanding the homomorphism count as
a sum over maps $\varphi:V(F)\to[n]$, the injective maps contribute
expectation
\[
 \frac{(n)_v}{n^v}t(F,W),
\]
while the number of non-injective maps is $O_F(n^{v-1})$.  Since every
summand lies in $[0,1]$, it follows that
\[
 \bbE\,t(F,\mathbb G(n,W))
 =t(F,W)+O_F(n^{-1}).
\]
In the variance expansion, two summands corresponding to maps with
disjoint images are independent.  The number of ordered pairs of maps
with intersecting images is $O_F(n^{2v-1})$, and each covariance has
absolute value at most one.  Hence
\[
 \operatorname{Var}\bigl(t(F,\mathbb G(n,W))\bigr)
 =O_F(n^{-1}).
\]
Chebyshev's inequality and a union bound over the fixed finite family
$F_1,\dots,F_q$ show that, for every sufficiently large $n$, there is a
deterministic realization $G_n$ for which all the required densities are
simultaneously arbitrarily close to their graphon values.  Choosing one
such realization for each $n$ proves the lemma.
\end{proof}

\section{The common \texorpdfstring{$K_r$}{Kr} marginal construction}\label{sec:common-law}

Let $G$ be a finite graph on $n$ vertices, let
\[
 p:=t(\K_2,G),
\]
and assume the strict range \eqref{eq:strict-range}.  By
\cref{thm:moon-moser}, $\Hom(\K_r,G)$ is nonempty.  Let
\[
 Y=(Y_1,\dots,Y_r)
\]
be uniformly distributed on $\Hom(\K_r,G)$.  The law of $Y$ is
exchangeable: for every permutation $\sigma\in S_r$,
\[
 (Y_{\sigma(1)},\dots,Y_{\sigma(r)})
 \stackrel{d}{=}
 (Y_1,\dots,Y_r).
\]
For $0\leq k\leq r$, define
\begin{equation}
 \cH_k:=H(Y_1,\dots,Y_k),
 \qquad
 \cH_0:=0,\label{eq:def-Hk}
\end{equation}
and, for $1\leq j\leq r$, define
\begin{equation}
 \alpha_j:=\cH_j-\cH_{j-1}
 =H(Y_j\mid Y_1,\dots,Y_{j-1}),
 \qquad
 \delta_j:=\log n-\alpha_j.\label{eq:def-alpha-delta}
\end{equation}
Since $Y_j$ takes at most $n$ values, $\alpha_j\leq\log n$, so
\begin{equation}
 \delta_j\geq0.\label{eq:delta-nonnegative}
\end{equation}

The next proposition is the point at which purity is relaxed.

\begin{proposition}[Entropy bound from one common clique law]
\label{prop:common-law-bound}
With the notation above,
\begin{equation}
 \log\homnum(H,G)
 \geq
 \sum_{i=1}^m\cH_{r_i}
 -
 \sum_{e\in E(T)}\cH_{s_e}
 =
 \sum_{j=1}^r b_j\alpha_j.\label{eq:common-law-hom-bound}
\end{equation}
Consequently,
\begin{equation}
 \log t(H,G)
 \geq
 -\sum_{j=1}^r b_j\delta_j.\label{eq:basic-deficit-bound}
\end{equation}
\end{proposition}

\begin{proof}
For each maximal clique $C_i$, choose an arbitrary injection
\[
 \iota_i:C_i\hookrightarrow[r].
\]
Define a probability measure $\mu_i$ on $\Hom(H[C_i],G)$ by taking the
law of the random assignment
\[
 v\longmapsto Y_{\iota_i(v)},
 \qquad v\in C_i.
\]
This assignment is a clique homomorphism because it is the restriction
of the homomorphism $Y:K_r\to G$.

We verify consistency on separators.  Let $e=ik\in E(T)$ and put
$S_e=C_i\cap C_k$.  The two restrictions
\[
 \iota_i|_{S_e},\ \iota_k|_{S_e}:S_e\hookrightarrow[r]
\]
are injections.  There is a permutation $\sigma\in S_r$ satisfying
\[
 \sigma(\iota_i(v))=\iota_k(v)
 \qquad(v\in S_e).
\]
Exchangeability of $Y$ therefore gives
\[
 \bigl(Y_{\iota_i(v)}\bigr)_{v\in S_e}
 \stackrel{d}{=}
 \bigl(Y_{\iota_k(v)}\bigr)_{v\in S_e}.
\]
Thus the marginals of $\mu_i$ and $\mu_k$ on $S_e$ agree.  Moreover, by
exchangeability, the entropy of every $k$-coordinate marginal of $Y$ is
$\cH_k$.  Hence
\[
 H(\mu_i)=\cH_{r_i},
 \qquad
 H(\nu_e)=\cH_{s_e}.
\]
Applying \cref{lem:junction-entropy} proves the first inequality in
\eqref{eq:common-law-hom-bound}.

Using $\cH_k=\sum_{j=1}^k\alpha_j$ and the definition
\eqref{eq:def-b}, we obtain
\[
 \begin{aligned}
 \sum_i\cH_{r_i}-\sum_e\cH_{s_e}
 &=\sum_i\sum_{j=1}^{r_i}\alpha_j
   -\sum_e\sum_{j=1}^{s_e}\alpha_j\\
 &=\sum_{j=1}^r
   \left(\#\{i:r_i\geq j\}-\#\{e:s_e\geq j\}\right)\alpha_j\\
 &=\sum_{j=1}^r b_j\alpha_j,
 \end{aligned}
\]
which completes \eqref{eq:common-law-hom-bound}.

Finally, by \eqref{eq:v-sum-b} and
$\alpha_j=\log n-\delta_j$,
\[
 \begin{aligned}
 \log t(H,G)
 &=\log\homnum(H,G)-v(H)\log n\\
 &\geq\sum_{j=1}^r b_j\alpha_j
       -\left(\sum_{j=1}^r b_j\right)\log n\\
 &=-\sum_{j=1}^r b_j\delta_j.
 \end{aligned}
\]
This proves \eqref{eq:basic-deficit-bound}.
\end{proof}

The Moon--Moser inequalities control every suffix sum of the deficits.

\begin{lemma}[Suffix bounds]\label{lem:suffix-bounds}
Let
\[
 h_k:=\homnum(\K_k,G),
 \qquad
 t_k:=t(\K_k,G),
 \qquad
 h_0:=1,
 \qquad
 t_0=t_1:=1.
\]
For $0\leq k<r$, define
\[
 D_k:=\sum_{j=k+1}^r\delta_j.
\]
Then
\begin{equation}
 D_k
 \leq
 -\log\frac{t_r}{t_k}
 \leq
 -\log\frac{A_r(p)}{A_k(p)}.\label{eq:suffix-bound-A}
\end{equation}
Equivalently, if
\begin{equation}
 \lambda_j:=-\log c_j(p),
 \qquad 1\leq j\leq r,\label{eq:def-lambda}
\end{equation}
then
\begin{equation}
 D_k\leq\sum_{j=k+1}^r\lambda_j.\label{eq:suffix-bound-lambda}
\end{equation}
\end{lemma}

\begin{proof}
Because $Y$ is uniform on $\Hom(K_r,G)$,
\[
 \cH_r=\log h_r.
\]
The random tuple $(Y_1,\dots,Y_k)$ is supported on
$\Hom(K_k,G)$, so \eqref{eq:support-entropy} gives
\[
 \cH_k\leq\log h_k.
\]
Using \eqref{eq:def-alpha-delta}, we therefore obtain
\[
 \begin{aligned}
 D_k
 &=\sum_{j=k+1}^r(\log n-\alpha_j)\\
 &=(r-k)\log n-(\cH_r-\cH_k)\\
 &\leq(r-k)\log n-\log h_r+\log h_k\\
 &=-\log\frac{h_r/n^r}{h_k/n^k}\\
 &=-\log\frac{t_r}{t_k}.
 \end{aligned}
\]
The second inequality in \eqref{eq:suffix-bound-A} follows from
\eqref{eq:moon-moser-iterated}.  Finally,
\[
 -\log\frac{A_r(p)}{A_k(p)}
 =-\sum_{j=k+1}^r\log c_j(p)
 =\sum_{j=k+1}^r\lambda_j,
\]
which proves \eqref{eq:suffix-bound-lambda}.
\end{proof}

For the extension to all chordal graphs in \cref{sec:all-chordal}, we
also need monotonicity of the deficits.

\begin{lemma}[Monotonicity of the entropy deficits]
\label{lem:delta-monotone}
The sequence $(\alpha_j)_{j=1}^r$ is nonincreasing, and hence
\begin{equation}
 0\leq\delta_1\leq\delta_2\leq\cdots\leq\delta_r.\label{eq:delta-monotone}
\end{equation}
\end{lemma}

\begin{proof}
For $1\leq j<r$, conditioning reduces entropy, so
\[
 \begin{aligned}
 \alpha_{j+1}
 &=H(Y_{j+1}\mid Y_1,\dots,Y_j)\\
 &\leq H(Y_{j+1}\mid Y_2,\dots,Y_j).
 \end{aligned}
\]
By exchangeability, the joint law of
$(Y_{j+1},Y_2,\dots,Y_j)$ is the same, up to a relabeling of coordinates,
as that of $(Y_j,Y_1,\dots,Y_{j-1})$.  Therefore
\[
 H(Y_{j+1}\mid Y_2,\dots,Y_j)
 =H(Y_j\mid Y_1,\dots,Y_{j-1})
 =\alpha_j.
\]
Thus $\alpha_{j+1}\leq\alpha_j$.  Since
$\delta_j=\log n-\alpha_j$, the sequence $(\delta_j)$ is nondecreasing.
Its nonnegativity was established in \eqref{eq:delta-nonnegative}.
\end{proof}

\section{Proof of the main theorem}\label{sec:main-proof}

We first prove the finite-graph statement in the strict range.

\begin{proposition}[Finite-graph form]\label{prop:finite-main}
Let $H$ satisfy the assumptions of \cref{thm:main}, and let $G$ be a
finite graph with
\[
 p=t(K_2,G)>1-\frac{1}{r-1}.
\]
Then
\[
 t(H,G)\geq\prod_{j=1}^r c_j(p)^{b_j}.
\]
\end{proposition}

\begin{proof}
Keep the notation $D_k$ and $\lambda_j$ from
\cref{lem:suffix-bounds}.  Abel summation, in the form obtained by
expanding each coefficient $b_j$ from its successive differences, gives
\begin{equation}
 \sum_{j=1}^r b_j\delta_j
 =
 b_1D_0
 +
 \sum_{k=1}^{r-1}(b_{k+1}-b_k)D_k.\label{eq:abel-b-delta}
\end{equation}
Indeed, the coefficient of $\delta_j$ on the right is
\[
 b_1+\sum_{k=1}^{j-1}(b_{k+1}-b_k)=b_j.
\]
By \cref{lem:b-properties}, $b_1\geq1$, and by the hypothesis
\eqref{eq:b-monotone}, every difference $b_{k+1}-b_k$ is nonnegative.
We may therefore apply the suffix estimates
\eqref{eq:suffix-bound-lambda} term by term in
\eqref{eq:abel-b-delta}.  This yields
\[
 \begin{aligned}
 \sum_{j=1}^r b_j\delta_j
 &\leq
 b_1\sum_{j=1}^r\lambda_j
 +
 \sum_{k=1}^{r-1}(b_{k+1}-b_k)
 \sum_{j=k+1}^r\lambda_j\\
 &=\sum_{j=1}^r b_j\lambda_j.
 \end{aligned}
\]
Combining this with \eqref{eq:basic-deficit-bound}, we obtain
\[
 \begin{aligned}
 \log t(H,G)
 &\geq-\sum_{j=1}^r b_j\delta_j\\
 &\geq-\sum_{j=1}^r b_j\lambda_j\\
 &=\sum_{j=1}^r b_j\log c_j(p)\\
 &=\log\prod_{j=1}^r c_j(p)^{b_j}.
\end{aligned}
\]
Exponentiating proves the proposition.
\end{proof}

\begin{proof}[Proof of \cref{thm:main}]
Let $W$ be a graphon with edge density $p$.

First suppose
\[
 1-\frac{1}{r-1}<p<1.
\]
Apply \cref{lem:finite-approximation} to the finite family consisting of
$H,K_2,\dots,K_r$.  We obtain finite graphs $G_n$ such that
\[
 t(H,G_n)\to t(H,W),
 \qquad
 p_n:=t(K_2,G_n)\to p.
\]
For all sufficiently large $n$,
$p_n>1-1/(r-1)$, so \cref{prop:finite-main} gives
\[
 t(H,G_n)\geq\prod_{j=1}^r c_j(p_n)^{b_j}.
\]
Letting $n\to\infty$ and using continuity of the right-hand side gives
\[
 t(H,W)\geq\prod_{j=1}^r c_j(p)^{b_j}=\PhiH(p),
\]
where the last equality is \eqref{eq:Phi-product}.

If $p=1-1/(r-1)$, then $c_r(p)=0$.  Since
$b_r\geq1$ by \cref{lem:b-properties}, the right-hand side of
\eqref{eq:main-bound} is zero, and the inequality follows from
$t(H,W)\geq0$.

Finally, if $p=1$, then $0\leq W\leq1$ and
$\int W=1$ imply that $W=1$ almost everywhere.  Hence
$t(H,W)=1=\PhiH(1)$.  This covers all cases.
\end{proof}

\begin{proof}[Proof of \cref{cor:pure}]
Suppose every maximal clique has size $r$.  Let
\[
 m_j:=\#\{i:r_i=j\},
 \qquad
 a_j:=\#\{e\in E(T):s_e=j\}.
\]
For $1\leq j<r$, direct subtraction from \eqref{eq:def-b} gives
\begin{equation}
 b_{j+1}-b_j=a_j-m_j.\label{eq:b-difference}
\end{equation}
In the pure case, $m_j=0$ for every $j<r$, so
$b_{j+1}-b_j=a_j\geq0$.  Thus \eqref{eq:b-monotone} holds, and the
corollary follows from \cref{thm:main}.
\end{proof}

\section{Structural interpretation and non-pure examples}
\label{sec:structure}

Equation \eqref{eq:b-difference} holds for every connected chordal graph,
not only in the pure case.  Therefore \eqref{eq:b-monotone} has the
following direct interpretation.

\begin{corollary}[Equivalent clique-tree condition]\label{cor:structural}
For $1\leq j<r$, let $m_j$ be the number of maximal cliques of size
exactly $j$, and let $a_j$ be the number of clique-tree edges whose
separator has size exactly $j$.  Then
\[
 b_1\leq b_2\leq\cdots\leq b_r
\]
if and only if
\begin{equation}
 m_j\leq a_j
 \qquad(1\leq j<r).\label{eq:structural-condition}
\end{equation}
The numbers $a_j$ are independent of the chosen clique tree.
\end{corollary}

\begin{proof}
The equivalence follows immediately from
$b_{j+1}-b_j=a_j-m_j$.  Since the numbers $m_j$ are intrinsic to $H$ and
the differences $b_{j+1}-b_j$ are intrinsic by
\cref{lem:b-properties}, each $a_j$ is also independent of the chosen
clique tree.
\end{proof}

Thus maximal cliques smaller than $r$ are allowed: at each size $j<r$,
their number must be at most the number of separators of the same size.

\begin{example}[A genuinely non-pure graph]\label{ex:nonpure}
Let $H$ have maximal cliques
\[
 C_1=\{a,b,c\},
 \qquad
 C_2=\{a,b,d\},
 \qquad
 C_3=\{c,e\}.
\]
A clique tree has edges $C_1C_2$ and $C_1C_3$, with separator sizes $2$
and $1$, respectively.  The maximal clique sizes are $3,3,2$, so
\[
 (b_1,b_2,b_3)=(1,2,2).
\]
The graph is not pure, but \cref{thm:main} applies and gives
\[
 t(H,W)\geq p^2(2p-1)^2,
 \qquad p\geq\frac12.
\]
Indeed,
\[
 \chi_H(x)=x(x-1)^2(x-2)^2.
\]
\end{example}

For clique number three, \eqref{eq:structural-condition} has a particularly
simple form.  Since $b_1=1$ and $b_2\geq1$, the only nonautomatic
condition is
\[
 \#\{\text{maximal edges}\}
 \leq
 \#\{\text{clique-tree separators that are edges}\}.
\]
Every chordal graph of clique number three satisfying this inequality
obeys the chromatic-polynomial lower bound for $p\geq1/2$.

\section{A bound for every connected chordal graph}
\label{sec:all-chordal}

When $(b_j)$ is not nondecreasing, the suffix estimates from
\cref{lem:suffix-bounds} cannot be combined directly with
\eqref{eq:basic-deficit-bound}: the coefficients
$b_{k+1}-b_k$ in \eqref{eq:abel-b-delta} may be negative.  The additional
fact \eqref{eq:delta-monotone} permits a canonical weakening.

\subsection{Pooling decreasing adjacent blocks}

We define a sequence $\bar b_1,\dots,\bar b_r$ by the following fixed
pool-adjacent-violators procedure.

Start with the ordered singleton blocks
\[
 \{1\},\{2\},\dots,\{r\},
\]
where a block $I$ is assigned the mean
\[
 \operatorname{av}(I):=\frac{1}{|I|}\sum_{j\in I}b_j.
\]
As long as two consecutive blocks $I,J$ satisfy
$\operatorname{av}(I)>\operatorname{av}(J)$, merge the leftmost such pair
into the single block $I\cup J$.  The number of blocks decreases at every
step, so the procedure terminates.  On each final block $I$, set
\begin{equation}
 \bar b_j:=\operatorname{av}(I)
 \qquad(j\in I).\label{eq:def-bar-b}
\end{equation}
By construction,
\[
 \bar b_1\leq\bar b_2\leq\cdots\leq\bar b_r.
\]

\begin{lemma}[Prefix and dual inequalities]\label{lem:pava}
Let
\[
 B_k:=\sum_{j=1}^k b_j,
 \qquad
 \bar B_k:=\sum_{j=1}^k\bar b_j.
\]
Then
\begin{equation}
 \bar B_k\leq B_k
 \quad(1\leq k<r),
 \qquad
 \bar B_r=B_r.\label{eq:pava-prefix}
\end{equation}
Consequently, for every nondecreasing real sequence
$z_1\leq\cdots\leq z_r$,
\begin{equation}
 \sum_{j=1}^r b_jz_j
 \leq
 \sum_{j=1}^r\bar b_jz_j.\label{eq:pava-dual}
\end{equation}
Moreover, every $\bar b_j$ is positive.
\end{lemma}

\begin{proof}
Initially the current sequence equals $(b_j)$, so all current prefix sums
are equal to $B_k$.  We show that one merge can only decrease the current
prefix sums and leaves the total sum unchanged.

Suppose two consecutive constant blocks
\[
 I=\{u,\dots,v\},
 \qquad
 J=\{v+1,\dots,w\}
\]
have lengths $\ell$ and $m$ and means $x>y$.  Their merged mean is
\[
 z:=\frac{\ell x+my}{\ell+m},
\]
so $y<z<x$ and
\begin{equation}
 \ell(x-z)=m(z-y).\label{eq:merge-balance}
\end{equation}
For a prefix ending at $k\in I$, replacing $x$ by $z$ changes the prefix
sum by
\[
 (k-u+1)(z-x)\leq0.
\]
For a prefix ending at $v+t\in J$, where $1\leq t\leq m$, the change is
\[
 \ell(z-x)+t(z-y)
 =-(m-t)(z-y)
 \leq0,
\]
where \eqref{eq:merge-balance} was used.  At the end of $J$ the change is
zero, and all later prefix sums are unchanged.  Induction over all merges
proves \eqref{eq:pava-prefix}.

Set
\[
 d_j:=\bar b_j-b_j,
 \qquad
 D_k:=\sum_{j=1}^k d_j=\bar B_k-B_k.
\]
Then $D_k\leq0$ for $k<r$ and $D_r=0$.  Abel summation gives
\[
 \begin{aligned}
 \sum_{j=1}^r d_jz_j
 &=\sum_{k=1}^{r-1}D_k(z_k-z_{k+1})\\
 &=-\sum_{k=1}^{r-1}D_k(z_{k+1}-z_k)\\
 &\geq0,
 \end{aligned}
\]
because both $-D_k$ and $z_{k+1}-z_k$ are nonnegative.  This is exactly
\eqref{eq:pava-dual}.  Finally, every $b_j$ is positive by
\cref{lem:b-properties}, and each $\bar b_j$ is an average of a nonempty
block of the $b_j$; hence $\bar b_j>0$.
\end{proof}

Define, for $p\in[0,1]$,
\begin{equation}
 \PsiH(p)
 :=
 \prod_{j=1}^r c_j(p)^{\bar b_j}.\label{eq:def-Psi}
\end{equation}
The exponents $\bar b_j$ may be rational rather than integral.

\begin{theorem}[Explicit bound for arbitrary chordal graphs]
\label{thm:all-chordal}
Let $H$ be any connected chordal graph with clique number $r\geq3$, and
let $\bar b$ be defined by \eqref{eq:def-bar-b}.  Then every graphon $W$
with edge density
\[
 p\geq1-\frac{1}{r-1}
\]
satisfies
\begin{equation}
 t(H,W)\geq\PsiH(p)
 =\prod_{j=1}^r c_j(p)^{\bar b_j}.\label{eq:all-chordal-bound}
\end{equation}
If $(b_j)$ is nondecreasing, then $\bar b=b$, and
\eqref{eq:all-chordal-bound} is exactly \cref{thm:main}.  In general,
\begin{equation}
 \PsiH(p)\leq\PhiH(p)
 \qquad
 \left(1-\frac{1}{r-1}\leq p\leq1\right).\label{eq:Psi-le-Phi}
\end{equation}
\end{theorem}

\begin{proof}
We first treat a finite graph $G$ in the strict range
$p>1-1/(r-1)$.  By \eqref{eq:basic-deficit-bound},
\[
 \log t(H,G)\geq-\sum_{j=1}^r b_j\delta_j.
\]
The deficits are nondecreasing by \cref{lem:delta-monotone}; therefore
\eqref{eq:pava-dual}, applied with $z_j=\delta_j$, gives
\[
 \sum_{j=1}^r b_j\delta_j
 \leq
 \sum_{j=1}^r\bar b_j\delta_j.
\]
Hence
\begin{equation}
 \log t(H,G)
 \geq
 -\sum_{j=1}^r\bar b_j\delta_j.\label{eq:bar-basic}
\end{equation}
Since $(\bar b_j)$ is nondecreasing and positive, the same Abel-summation
argument as in \cref{prop:finite-main}, together with the suffix bounds
from \cref{lem:suffix-bounds}, gives
\[
 \sum_{j=1}^r\bar b_j\delta_j
 \leq
 \sum_{j=1}^r\bar b_j\lambda_j.
\]
Substituting this into \eqref{eq:bar-basic} yields
\[
 \log t(H,G)
 \geq
 -\sum_{j=1}^r\bar b_j\lambda_j
 =
 \log\prod_{j=1}^r c_j(p)^{\bar b_j}.
\]
This proves the finite-graph inequality in the strict range.

The passage to graphons, the boundary
$p=1-1/(r-1)$, and the case $p=1$ are handled exactly as in the proof of
\cref{thm:main}.  At the lower boundary, $c_r(p)=0$ and
$\bar b_r>0$, so the right-hand side of
\eqref{eq:all-chordal-bound} is zero.

If $(b_j)$ is nondecreasing, the pooling procedure performs no merge, so
$\bar b=b$.  Finally, in the strict range, the sequence
$\lambda_j=-\log c_j(p)$ is nondecreasing.  Applying
\eqref{eq:pava-dual} with $z_j=\lambda_j$ gives
\[
 \sum_j b_j\lambda_j
 \leq
 \sum_j\bar b_j\lambda_j.
\]
After multiplying by $-1$ and exponentiating, this is
$\PsiH(p)\leq\PhiH(p)$.  The two boundary cases follow by continuity.
\end{proof}

\begin{example}[A triangle with pendant leaves]\label{ex:leaves}
Let $H_L$ be obtained from a triangle by attaching $L$ leaves to each of
its three vertices.  Its maximal cliques consist of one triangle and
$3L$ pendant edges.  In a clique tree, every pendant-edge clique is
joined to the triangle clique along a one-vertex separator.  Therefore
\[
 (b_1,b_2,b_3)=(1,3L+1,1).
\]
For $L\geq1$, the pooling procedure merges the final two entries and gives
\[
 (\bar b_1,\bar b_2,\bar b_3)
 =
 \left(1,\frac{3L+2}{2},\frac{3L+2}{2}\right).
\]
Thus \cref{thm:all-chordal} yields
\begin{equation}
 t(H_L,W)
 \geq
 \bigl[p(2p-1)\bigr]^{(3L+2)/2},
 \qquad p\geq\frac12.\label{eq:leaf-bound}
\end{equation}
By contrast,
\[
 \Phi_{H_L}(p)=p^{3L+1}(2p-1).
\]
The downward jump from $b_2=3L+1$ to $b_3=1$ is exactly the obstruction
to applying \cref{thm:main}; the pooled bound
\eqref{eq:leaf-bound} records the entropy loss forced by that jump.
\end{example}

\section{Conclusion}

The common source of all local clique laws is the decisive modification.
For a pure chordal graph, every bag already carries the full uniform
$K_r$-law.  For a non-pure graph, a bag of size $r_i<r$ instead carries
an $r_i$-coordinate marginal of the same uniform $K_r$-law.  This restores
separator consistency automatically.  The resulting entropy deficits
are not controlled coordinate by coordinate; Moon--Moser controls their
suffix sums.  The condition $b_1\leq\cdots\leq b_r$ is precisely what
allows those suffix inequalities to be combined with nonnegative
coefficients.  When the condition fails, monotonicity of the entropy
deficits and adjacent-block pooling give the weaker universal bound of
\cref{thm:all-chordal}.

\begin{thebibliography}{9}

\bibitem{Agnarsson}
G.~Agnarsson,
\emph{On chordal graphs and their chromatic polynomials},
Math. Scand. \textbf{93} (2003), 240--246.

\bibitem{Lovasz}
L.~Lov\'asz,
\emph{Large Networks and Graph Limits},
American Mathematical Society, Providence, RI, 2012.

\bibitem{MoonMoser}
J.~W. Moon and L.~Moser,
\emph{On a problem of Tur\'an},
Magyar Tud. Akad. Mat. Kutat\'o Int. K\"ozl. \textbf{7} (1962), 283--286.

\bibitem{Zhao}
Y.~Zhao,
\emph{Graph Theory and Additive Combinatorics: Exploring Structure and Randomness},
Cambridge University Press, Cambridge, 2023.

\end{thebibliography}

\end{document}
